\documentclass[../main]{subfiles}
\begin{document}
\chapter{Transformaciones}
\img{images/05/modulogas.jpg}{Moléculas de agua en estado gaseoso. }
\info{
Contenido}{
  Transformaciones: isócora, isóbara, isotérmica, adiabática,
  isoentrópica, politrópica, problemas.
}

\info{Transformación termodinámica}{
  Una transformación termodinámica se refiere a los cambios que
  experimenta un sistema termodinámico cuando pasa de un estado de
  equilibrio a otro. En termodinámica, un sistema se compone de un
  conjunto de partículas o sustancias que interactúan entre sí y con
  su entorno, y una transformación termodinámica describe cómo
  cambian las propiedades del sistema durante ese proceso.
}
\begin{center}
  \begin{tabular}{l|c|c}
    Gas & Formula Química & $R \left[\dfrac{kgm}{kg ^\circ K}\right]$  \\ \hline
    Acetileno & $C_2 H_2$ & 32,56 \\
    Ácido sulfhídrico & $SH_2$ & 24,87 \\
    Anhídrido carbónico & $CO_2$& 19,26 \\
    Anhídrido sulfuroso& $SO_2$&13,23 \\
    Aire&...&29,26 \\
    Amoníaco& $NH_3$&49,77 \\
    Argón& A& 21,22\\
    Helio& $He$&211,79 \\
    Hidrógeno&$H_2$&420,55 \\
    Metano& $CH_4$ &52,85 \\
    Monóxido de carbono&$CO$&20,27 \\
    Nitrógeno&$N_2$& 30,26\\
    Oxígeno&$O_2$& 26,44\\
    Vapor de agua& $H_2 O$& 47,05\\
  \end{tabular}
\end{center}

\actividad{ Actividad: Contestar usando el libro}
\begin{enumerate}
  \item Explicar transformación cerrada y abierta.
  \item Explicar transformación reversible e irreversible.
  \item Explicar y realizar el diagrama de Clapeyron, formulas
    principales (recuadradas en el libro) en sistemas no elástico las
    transformaciones:
    \begin{enumerate}
      \item Isócora.
      \item Isóbara.
      \item Isotermica.
      \item Adiabática.
      \item Isoentrópica.
      \item Politrópica
    \end{enumerate}
\end{enumerate}
\actividad{ Ejercicios: Resolver}
\begin{enumerate}
  \item Un recipiente rígido contiene $10kg$ de aire a $27^\circ C$ y
    recibe del medio exterior una cantidad de calor tal, que su
    presión pasa de $P_1=1 atm$ a $P_2=2 atm$. Determinar:
    \begin{enumerate}
      \item El estado inicial y el estado final del aire.
      \item La cantidad de calor que recibió desde el medio exterior.
      \item El valor del trabajo externo.
    \end{enumerate}
  \item Un kilogramo de oxigeno $O_2$ contenido en un recipiente de
    paredes rígidas, que se encuentra a la presión de \SI{3}{\atm} y a la
    temperatura de $47^\circ C$, entrega al medio exterior una
    cantidad de calor igual a 6kcal. Calcular:
    \begin{enumerate}
      \item La temperatura final y la presión final de oxígeno.
      \item La variación de energía interna que experimentó.
      \item El trabajo externo desarrollado.
    \end{enumerate}
  \item Un recipiente rígido, de paredes adiabáticas, contiene
    $0,5m^2$ de ácido sulfhídrico a una presión de $12 atm$. Mediante
    paletas colocadas en su interior y movidas por un motor, se agita
    el gas hasta que la presión se eleve a $13atm$. Calcular:
    \begin{enumerate}
      \item La variación de energía interna que experimentó la masa gaseosa.
      \item El trabajo desarrollado.
    \end{enumerate}
  \item Un globo que contiene 100kg de hidrógeno a $27^\circ C$, está
    sometido a la presión atmosférica normal. Suponiendo que la
    temperatura del gas descienda a $23^\circ C$ Calcular:
    \begin{enumerate}
      \item El estado inicial y el estado final del hidrógeno.
      \item el calor desarrollado.
      \item el trabajo externo.
    \end{enumerate}
  \item Una masa gaseosa de $5kg$ de peso, se encuentra a la
    temperatura de $47^\circ C$ y al entregarle $120kcal$ sufre una
    "transformación isobárica". Calcular la temperatura del estado
    final del gas.
  \item Calcular la cantidad de calor y el trabajo externo
    desarrollado cuando $1m^2$ de helio a $47^\circ C$ se comprime
    isotérmicamente desde la presión de $1 atm$ hasta una presión de $ 5 atm$
  \item Un metro cúbico de amoníaco a $27^\circ C $ y a la presión de
    $1atm$ es comprimido adiabáticamente hasta una presión de $6atm$.
    Determinar el estado inicial y el estado final del gas y calcular
    el trabajo externo.
  \item El Submarino \emph{“Titan” de OceanGate} era un artefacto
    cilíndrico de 6,80 metros de largo y 2,80 metros de diámetro
    lleno de 1,43kg de nitrógeno a presión de 1 atmósfera. Antes de
    la implosión catastrófica estuvo expuesto a una temperatura de
    $t_1=4^\circ C$ que es la que hay a \SI{3800}{\m} de profundidad con
    una presión de 370 atmósferas. Determinar el estado final del
    submarino después de la implosión adiabática.
  \item Calcular el trabajo externo y el calor transmitido cuando se
    comprime $1kg$ de nitrógeno de $17^\circ C$ y a la presión
    atmosférica normal, hasta una presión de $3.1atm$, si durante la
    compresión el volumen y la presión mantuvieron la siguiente
    relación: $p V^{1,2}=cte$.
  \item La expansión de 10kg de helio a la temperatura de $37^\circ
    C$ a la presión de $4 atm$, se produce a través de una
    "transformación politrópica", con exponente $1,1$ hasta que el
    gas alcanza la temperatura de $-13^\circ C$. Calcular:
    \begin{enumerate}
      \item El estado inicial y el estado final de la masa gaseosa.
      \item El trabajo externo y el calor desarrollado.
    \end{enumerate}
  \item Un compresor impele $200kg$ de amoníaco por hora, a la
    presión de $4 atm$, tomándolo a la presión de $1 atm$ y a la
    temperatura de $17^\circ C$. Suponiendo que la compresión sea una
    "transformación politrópica de exponente $1.2$. Calcular:
    \begin{enumerate}
      \item el estado inicial y final del amoníaco.
      \item el trabajo externo y el calor desarrollado.
      \item la potencia teórica necesaria para efectuar la compresión.
    \end{enumerate}
\end{enumerate}
\end{document}

